\ifdefined\gkver\else\def\gkver{student}\fi
\documentclass[\gkver]{gaokaozhenti}
\begin{document}

\chapter{2012年江苏}

\begin{enumerate}
\item
%% number: 1
%% paperName: 2012年江苏高考第 1 题 3 分
%% typeId: 1
%% type: 单选题
%% chapter: 电场
%% point: 电场强度
%% method: 无
%% score: 3
%% degree: 650
%% duplicateId: 0
%% body:
真空中，A、B两点与点电荷 $Q$ 的距离分别为 $r$ 和 3$r$，则A、B两点的电场强度大小之比为\xzanswer{C}

\fourchoices[answer=C]
{3$:$1}
{1$:$3}
{9$:$1}
{1$:$9}
%% answer: C
%% memo:
\memoanswer{
根据点电荷周围的电场强度表达式 $E=\frac{kQ}{r^{2}}$，可知点电荷周围的电场强度的大小与距离的二次方成反比，A、B 两点与点电荷 $Q$ 的距离之比为 1$:$3，所以电场强度大小之比为 9$:$1，C 项正确。
}

\item
%% number: 2
%% paperName: 2012年江苏高考第 2 题 3 分
%% typeId: 1
%% type: 单选题
%% chapter: 电路
%% point: 电容
%% method: 无
%% score: 3
%% degree: 650
%% duplicateId: 0
%% body:
一充电后的平行板电容器保持两极板的正对面积、间距和电荷量不变，在两极板间插入一电介质，其电容 $C$ 和两极板间的电势差 $U$ 的变化情况是\xzanswer{B}

\fourchoices[answer=B]
{$C$ 和 $U$ 均增大}
{$C$ 增大，$U$ 减小}
{$C$ 减小，$U$ 增大}
{$C$ 和 $U$ 均减小}
%% answer: B
%% memo:
\memoanswer{
由平行板电容器的电容的表达式 $C=\frac{\varepsilon S}{4\pi kd}$，当两极板之间插入一电介质，$\varepsilon$ 变大，则 $C$ 变大，由 $U=\frac{Q}{C}$ 可知，在电荷量 $Q$ 不变的情况下，两极板间的电势差 $U$ 将减小，B 项正确。
}

\item
%% number: 3
%% paperName: 2012年江苏高考第 3 题 3 分
%% typeId: 1
%% type: 单选题
%% chapter: 功和能
%% point: 功率
%% method: 无
%% score: 3
%% degree: 450
%% duplicateId: 0
%% body:
如图所示，细线的一端固定于O点，另一端系一小球。在水平拉力作用下，小球以恒定速率在竖直平面内由A点运动到B点。在此过程中拉力的瞬时功率变化情况是\xzanswer{A}
\onepicture[w=2.6cm, align=right]{figs/03.png}

\fourchoices[answer=A]
{逐渐增大}
{逐渐减小}
{先增大，后减小}
{先减小，后增大}
%% answer: A
%% memo:
\memoanswer{
设 $F$ 与速度 $v$ 的夹角为 $\theta$，则 $P=Fv\cos\theta$。小球以恒定速率运动，在切线上（速度方向上）合力为 0，即 $mg\sin\theta=F\cos\theta$，所以 $P=mg\sin\theta$，随 $\theta$ 增大，$P$ 增大。
}

\item
%% number: 4
%% paperName: 2012年江苏高考第 4 题 3 分
%% typeId: 1
%% type: 单选题
%% chapter: 运动和力
%% point: 变加速直线运动
%% method: 无
%% score: 3
%% degree: 450
%% duplicateId: 0
%% body:
一只皮球竖直向上抛出，皮球运动时受到空气阻力的大小与速度的大小成正比。下列描绘将皮球在上升过程中加速度大小 $a$ 与时间 $t$ 关系的图象，可能正确的是\xzanswer{C}

\fourchoices[answer=C, ispicture=true, h=2.2cm]
{figs/04a.png}
{figs/04b.png}
{figs/04c.png}
{figs/04d.png}
%% answer: C
%% memo:
\memoanswer{
由牛顿第二定律：$-(mg+kv)=ma$，加速度 $a=-(g+\frac{kv}{m})$，皮球在上升过程中做减速运动，随着 $v$ 的减小，$a$ 减小，但最后不等于 0。加速度越小，速度减小得越慢，所以选 C。
}

\item
%% number: 5
%% paperName: 2012年江苏高考第 5 题 3 分
%% typeId: 1
%% type: 单选题
%% chapter: 运动和力
%% point: 连接体
%% method: 整体与隔离
%% score: 3
%% degree: 450
%% duplicateId: 0
%% body:
如图所示，一夹子夹住木块，在力 $F$ 作用下向上提升。夹子和木块的质量分别为 $m$、$M$，夹子与木块两侧间的最大静摩擦力均为 $f$。若木块不滑动，力 $F$ 的最大值是\xzanswer{A}
\onepicture[w=2.4cm, align=right]{figs/05.png}

\fourchoices[answer=A]
{$\frac{2f(m+M)}{M}$}
{$\frac{2f(m+M)}{m}$}
{$\frac{2f(m+M)}{M}-(m+M)g$}
{$\frac{2f(m+M)}{m}+(m+M)g$}
%% answer: A
%% memo:
\memoanswer{
当木块所受的摩擦力最大时加速度最大，整体法 $F_{m}-(M+m)g=ma_{m}$，隔离法，对木块，$2f-Mg=Ma_{m}$，解得 $F_{m}=\frac{2f(m+M)}{M}$。
}

\item
%% number: 6
%% paperName: 2012年江苏高考第 6 题 4 分
%% typeId: 2
%% type: 多选题
%% chapter: 曲线运动
%% point: 平抛运动
%% method: 无
%% score: 4
%% degree: 450
%% duplicateId: 0
%% body:
如图所示，相距 $l$ 的两小球A、B位于同一高度 $h$（$l$，$h$均为定值）。将A向B水平抛出的同时，B自由下落。A、B与地面碰撞前后，水平分速度不变，竖直分速度大小不变、方向相反．不计空气阻力及小球与地面碰撞的时间，则\xzanswer{AD}
\onepicture[w=3.4cm, align=right]{figs/06.png}

\fourchoices[answer=AD]
{A、B在第一次落地前能否相碰，取决于A的初速度}
{A、B在第一次落地前若不碰，此后就不会相碰}
{A、B不可能运动到最高处相碰}
{A、B一定能相碰}
%% answer: AD
%% memo:
\memoanswer{
A 做平抛运动，竖直方向的分运动为自由落体运动，满足关系式 $h=\frac{1}{2}gt^{2}$，水平方向上为匀速直线运动，满足关系式 $x=vt$，B 做自由落体运动，因为 A、B 从同一高度开始运动，因此两者在空中同一时刻处于同一高度，即使两者与地面撞击，反弹后在空中也是同一时刻处于同一高度，而 A 在水平方向一直向右运动，因此 A、B 肯定会相碰，D 项正确；当 A 的水平速度 $v$ 足够大时，有可能在 B 未落地前二者相碰，因此 A、B 在第一次落地前能否相碰，取决于 A 的初速度，A 项正确。
}

\item
%% number: 7
%% paperName: 2012年江苏高考第 7 题 4 分
%% typeId: 2
%% type: 多选题
%% chapter: 电磁感应
%% point: 电磁感应现象
%% method: 无
%% score: 4
%% degree: 650
%% duplicateId: 0
%% body:
某同学设计的家庭电路保护装置如图所示，铁芯左侧线圈 L$_{1}$ 由火线和零线并行绕成．当右侧线圈 L$_{2}$ 中产生电流时，电流经放大器放大后，使电磁铁吸起铁质开关K，从而切断家庭电路．仅考虑 L$_{1}$，在铁芯中产生的磁场，下列说法正确的有\xzanswer{ABD}
\onepicture[w=4.6cm, align=right]{figs/07.png}

\fourchoices[answer=ABD]
{家庭电路正常工作时，L$_{2}$ 中的磁通量为零}
{家庭电路中使用的电器增多时，L$_{2}$ 中的磁通量不变}
{家庭电路发生短路时，开关K将被电磁铁吸起}
{地面上的人接触火线发生触电时，开关K将被电磁铁吸起}
%% answer: ABD
%% memo:
\memoanswer{
由于火线和零线并行绕制，所以在家庭电路正常工作时，火线和零线的电流大小相等，方向相反，因此合磁通量为零，L$_{2}$ 中的磁通量为零，A 项正确；当家庭电路中使用的电器增多时，火线和零线的电流仍然大小相等，方向相反，L$_{2}$ 中的磁通量不变，B 项正确；家庭电路发生短路时，火线和零线的电流同时增大，合磁通量仍然为零，因此开关 K 不会被电磁铁吸起，C 项错误；当地面上的人接触火线发生触电时，火线的电流突然变大，即 L$_{1}$ 中的磁场发生变化，导致 L$_{2}$ 中的磁通量变化，L$_{2}$ 中产生感应电流，电磁铁将开关 K 吸起，D 项正确。
}

\item
%% number: 8
%% paperName: 2012年江苏高考第 8 题 4 分
%% typeId: 2
%% type: 多选题
%% chapter: 曲线运动
%% point: 人造地球卫星
%% method: 无
%% score: 4
%% degree: 450
%% duplicateId: 0
%% body:
2011年8月，“嫦娥二号”成功进入了环绕“日地拉格朗日点”的轨道，我国成为世界上第三个造访该点的国家。如图所示，该拉格朗日点位于太阳和地球连线的延长线上，一飞行器处于该点，在几乎不消耗燃料的情况下与地球同步绕太阳做圆周运动。则此飞行器的\xzanswer{AB}
\onepicture[w=2.6cm, align=right]{figs/08.png}

\fourchoices[answer=AB]
{线速度大于地球的线速度}
{向心加速度大于地球的向心加速度}
{向心力仅有太阳的引力提供}
{向心力仅由地球的引力提供}
%% answer: AB
%% memo:
\memoanswer{
飞行器与地球绕太阳同步做圆周运动，它们的周期相同，角速度也相同，由 $v=\omega r$，$a=\omega^{2}r$，可知半径越大，线速度越大，向心加速度也越大，A、B 正确；飞行器的向心力由太阳和地球的合力来提供，C、D 错误。
}

\item
%% number: 9
%% paperName: 2012年江苏高考第 9 题 4 分
%% typeId: 2
%% type: 多选题
%% chapter: 磁场
%% point: 带电粒子在磁场中的运动
%% method: 无
%% score: 4
%% degree: 450
%% duplicateId: 0
%% body:
如图所示，MN是磁感应强度为 $B$ 的匀强磁场的边界。一质量为 $m$、电荷量为 $q$ 的粒子在纸面内从O点射入磁场。若粒子速度为 $v_{0}$，最远能落在边界上的A点。下列说法正确的有\xzanswer{BC}
\onepicture[w=3.4cm, align=right]{figs/09.png}

\fourchoices[answer=BC]
{若粒子落在A点的左侧，其速度一定小于 $v_{0}$}
{若粒子落在A点的右侧，其速度一定大于 $v_{0}$}
{若粒子落在A点左右两侧 $d$ 的范围内，其速度不可能小于 $v_{0}-\frac{qBd}{2m}$}
{若粒子落在A点左右两侧 $d$ 的范围内，其速度不可能大于 $v_{0}+\frac{qBd}{2m}$}
%% answer: BC
%% memo:
\memoanswer{
带电粒子沿垂直边界的方向射入磁场时，落在边界上的点离出发点最远，当入射方向不是垂直边界的方向时，落在边界上的点与出发点的距离将小于这个距离，即速度大于或等于 $v_{0}$，但入射方向不是 $90^{\circ}$ 时，粒子有可能落在 A 点的左侧，A 项错误；但粒子要落在 A 点的右侧，其速度一定要大于临界速度 $v_{0}$，B 项正确；设 OA 之间距离为 $L$，若粒子落在 A 点两侧 $d$ 范围内，则以最小速度 $v$ 入射的粒子做圆周运动的直径应为 $L-d$，由洛伦兹力提供向心力，$qvB=\frac{mv^{2}}{\frac{L-d}{2}}$，$qv_{0}B=\frac{mv_{0}^{2}}{\frac{L}{2}}$，解得 $v=v_{0}-\frac{qBd}{2m}$，C 项正确；由于题中没有强调粒子的入射方向，因此无法确定速度的最大值，D 项错误。
}

\item
%% number: 10
%% paperName: 2012年江苏高考第 10 题 8 分
%% typeId: 4
%% type: 实验题
%% chapter: 电路
%% point: 多用表的使用
%% method: 无
%% score: 8
%% degree: 650
%% duplicateId: 0
%% body:
如图1所示的黑箱中有二只完全相同的电学元件，小明使用多用电表对其进行探测。

\onepicture[w=8.4cm]{figs/10.png}

\begin{enumerate}
	\item
	在使用多用电表前，发现指针不在左边“0”刻度线处，应先调整图2中多用电表的\tkanswer[1.2cm]{A}（选填“A”、“B”或“C”）。
	\item
	在用多用电表的直流电压挡探测黑箱 $a$、$b$ 接点间是否存在电源时，一表笔接 $a$，另一表笔应\tkanswer[1.4cm]{短暂}（选填“短暂”或“持续”）接 $b$，同时观察指针偏转情况。
	\item
	在判定黑箱中无电源后，将选择开关旋至“$\times1$”挡，调节好多用电表，测量各接点间的阻值。测量中发现，每对接点间正反向阻值均相等，测量记录如下表。两表笔分别接 $a$、$b$ 时，多用电表的示数如图2所示。
	
	请将记录表补充完整，并在答题卡的黑箱图中画出一种可能的电路。
	
	\begin{center}
	\begin{tabular}{|c|c|}
	\hline
	两表笔接的接点 & 多用电表的示数 \\
	\hline
	a，b & \tkanswer[1.6cm]{5.0}$\UO$ \\
	\hline
	a，c & 10.0$\UO$ \\
	\hline
	b，c & 15.0$\UO$ \\
	\hline
	\end{tabular}
	\end{center}
\end{enumerate}
%% answer: 
\jdanswer{
\begin{enumerate}
	\item
	A
	\item
	短暂
	\item
	5.0（见图）
\end{enumerate}
}
%% memo:
\memoanswer{
（1）在使用多用电表时发现指针不在左边“0”刻度线处，则需要进行机械调零，即调节 A；
（2）在用直流电压挡探测黑箱 $a$、$b$ 接点间是否有电源时，一个表笔接 $a$，另一个表笔应短暂接 $b$，防止电流过大烧毁电表；
（3）每对接点间正反向电阻相等，说明黑箱中只有电阻，由于选择开关选择的是“$\times1$”挡，因此两表笔分别接 $a$、$b$ 时，多用电表的示数从表盘中读出为 $5.0\UO$；根据表格中数据，可知各接点间的电阻连接如图所示。

\onepicture[h=3.2cm, align=right]{figs/10_an.png}
}

\item
%% number: 11
%% paperName: 2012年江苏高考第 11 题 10 分
%% typeId: 4
%% type: 实验题
%% chapter: 功和能
%% point: 动能定理
%% method: 无
%% score: 10
%% degree: 450
%% duplicateId: 0
%% body:
为测定木块与桌面之间的动摩擦因数，小亮设计了如图所示的装置进行实验。实验中，当木块A位于水平桌面上的O点时，重物B刚好接触地面。将A拉到P点，待B稳定后静止释放，A最终滑到Q点。分别测量OP、OQ的长度 $h$ 和 $s$。改变 $h$，重复上述实验，分别记录几组实验数据。

\onepicture[w=7.2cm, align=right]{figs/11a.png}

\begin{enumerate}
	\item
	实验开始时，发现A释放后会撞到滑轮。请提出两个解决方法。
	\item
	请根据下表的实验数据作出 $s-h$ 关系的图象。
	
	\begin{center}
	\begin{tabular}{|c|c|c|c|c|c|}
	\hline
	$h$（cm） & 20.0 & 30.0 & 40.0 & 50.0 & 60.0 \\
	\hline
	$s$（cm） & 19.5 & 28.5 & 39.0 & 48.0 & 56.5 \\
	\hline
	\end{tabular}
	\end{center}
	
	\onepicture[w=5.2cm, align=right]{figs/11b.png}
	\item
	实验测得A、B的质量分别为 $m=0.4\Ukg$、$M=0.50\Ukg$。根据 $s-h$ 图象可计算出A木块与桌面间的动摩擦因数 $\mu=$\tkanswer[1.2cm]{0.4}。（结果保留一位有效数字）
	\item
	实验中，滑轮轴的摩擦会导致 $\mu$ 的测量结果\tkanswer[1.2cm]{偏大}（选填“偏大”或“偏小”）。
\end{enumerate}
%% answer: 
\jdanswer{
\begin{enumerate}
	\item
	减小 B 的质量；增加细线的长度；（或增大 A 的质量；降低 B 的起始高度）
	\item
	见下图
	\item
	0.4
	\item
	偏大
\end{enumerate}
}
%% memo:
\memoanswer{
（1）实验开始时发现 A 释放后会撞到滑轮，主要是加速度过大或加速时间过长，可以通过减小 B 的质量或增大 A 的质量来减小加速度，通过增加细线的长度或降低 B 的起始高度来缩短加速时间。
（2）$s-h$ 图象如图所示。
（3）本实验测动摩擦因数的原理是动能定理，即由 $Mgh-\mu mgh=\frac{1}{2}(M+m)v^{2}$，$-\mu mgs=-\frac{1}{2}mv^{2}$，求得 $s=\frac{M-\mu m}{\mu(M+m)}h$，图象的斜率 $k=\frac{M-\mu m}{\mu(M+m)}$，即 $\frac{0.5-0.4\mu}{\mu(0.5+0.4)}=\frac{56}{60}$，解得 $\mu=0.4$。
（4）本实验测动摩擦因数的原理是动能定理，如果考虑克服滑轮摩擦做功 $W$，则 $Mgh-\mu mgh-W=\frac{1}{2}(M+m)v^{2}$，$-\mu mgs=-\frac{1}{2}mv^{2}$，求得 $\mu=\frac{Mgh-W}{mgh+(M+m)gs}$，如果忽略克服滑轮摩擦做功，则动摩擦因数偏大。

\onepicture[h=4.2cm, align=right]{figs/11_an.png}
}

\item
%% number: 12
%% paperName: 2012年江苏高考第 12 题 4 分
%% typeId: 6
%% type: 计算题
%% chapter: 光学
%% point: 光电效应
%% method: 无
%% score: 4
%% degree: 650
%% duplicateId: 0
%% body:
【选做题】本题包括 A、B、C 三小题，请选定其中两小题，并在相应的答题区域内作答。若多做，则按 A、B 两小题评分。

\textbf{A．}（选修模块 3-3）

\begin{enumerate}
	\item
	下列现象中，能说明液体存在表面张力的有\xzanswer{AB}
	
	\fourchoices[answer=AB]
	{水黾可以停在水面上}
	{叶面上的露珠呈球形}
	{滴入水中的红墨水很快散开}
	{悬浮在水中的花粉做无规则运动}
	\item
	密闭在钢瓶中的理想气体，温度升高时压强增大。从分子动理论的角度分析，这是由于分子热运动的\tkanswer[1.4cm]{平均动能}增大了。该气体在温度 $T_{1}$、$T_{2}$ 时的分子速率分布图象如图所示，则 $T_{1}$\tkanswer[1.0cm]{小于}（选填“大于”或“小于”）$T_{2}$。
	
	\onepicture[w=6.0cm, align=right]{figs/12a.png}
	\item
	如图所示，一定质量的理想气体从状态A经等压过程到状态B。此过程中，气体压强 $p=1.0\times10^{5}\UPa$，吸收的热量 $Q=7.0\times10^{2}\UJ$，求此过程中气体内能的增量。
	
	\onepicture[w=5.2cm, align=right]{figs/12b.png}
\end{enumerate}

\textbf{B．}（选修模块 3-4）

\begin{enumerate}
	\item
	如图所示，白炽灯的右侧依次平行放置偏振片P和Q，A点位于P、Q之间，B点位于Q右侧。旋转偏振片P，A、B两点光的强度变化情况是\xzanswer{C}
	
	\onepicture[w=5.2cm]{figs/12c.png}
	
	\fourchoices[answer=C]
	{A、B均不变}
	{A、B均有变化}
	{A不变，B有变化}
	{A有变化，B不变}
	\item
	“测定玻璃的折射率”实验中，在玻璃砖的一侧竖直插两个大头针A、B，在另一侧再竖直插两个大头针C、D。在插入第四个大头针D时，要使它\tkanswer[2.4cm]{挡住C及A、B的像}。如图是在自纸上留下的实验痕迹，其中直线 a、a$'$ 是描在纸上的玻璃砖的两个边。根据该图可算得玻璃的折射率 $n=$\tkanswer[1.2cm]{1.8}。（计算结果保留两位有效数字）
	
	\onepicture[w=3.6cm, align=right]{figs/12d.png}
	\item
	地震时，震源会同时产生两种波，一种是传播速度约为 $3.5\Ukms$ 的S波，另一种是传播速度约为 $7.0\Ukms$ 的P波。一次地震发生时，某地震监测点记录到首次到达的P波比首次到达的S波早 $3\Umin$。假定地震波沿直线传播，震源的振动周期为 $1.2\Us$，求震源与监测点之间的距离 $x$ 和S波的波长 $\lambda$。
\end{enumerate}

\textbf{C．}（选修模块 3-5）

\begin{enumerate}
	\item
	如图所示是某原子的能级图，a、b、c 为原子跃迁所发出的三种波长的光。在下列该原子光谱的各选项中，谱线从左向右的波长依次增大，则正确的是\xzanswer{C}
	
	\onepicture[w=3.0cm]{figs/12e.png}
	
	\fourchoices[answer=C, ispicture=true, h=2.4cm]
	{figs/12e1.png}
	{figs/12e2.png}
	{figs/12e3.png}
	{figs/12e4.png}
	\item
	一个中子与某原子核发生核反应，生成一个氘核，其核反应方程式为\tkanswer[3.8cm]{}。该反应放出的能量为 $Q$，则氘核的比结合能为\tkanswer[1.6cm]{$\frac{Q}{2}$}。
	\item
	A、B两种光子的能量之比为 2$:$1，它们都能使某种金属发生光电效应，且所产生的光电子最大初动能分别为 $E_{A}$、$E_{B}$。求A、B两种光子的动量之比和该金属的逸出功。
\end{enumerate}
%% answer: 
\jdanswer{
\begin{enumerate}
	\item
	A：（1）AB；（2）平均动能，小于；（3）$\Delta U=5.0\times10^{2}\UJ$。
	\item
	B：（1）C；（2）挡住C及A、B的像，1.8（1.6$\sim$1.9 都算对）；（3）$x=1260\Ukm$，$\lambda=4.2\Ukm$。
	\item
	C：（1）C；（2）\ce{^{1}_{0}n + ^{1}_{1}H -> ^{2}_{1}H}，$\frac{Q}{2}$；（3）2$:$1，$E_{A}-2E_{B}$。
\end{enumerate}
}
%% memo:
\memoanswer{
\textbf{A．}（1）滴入水中的红墨水很快散开是扩散现象，悬浮在水中的花粉的运动是布朗运动，它说明水分子永不停息地做无规则运动，故能说明液体存在表面张力的是 A、B。
（2）温度升高，气体分子的平均动能增加，随着温度的升高，分子速率分布曲线的峰值向分子速率增大的方向移动，因此 $T_{1}$ 小于 $T_{2}$。
（3）等压变化 $\frac{V_{A}}{T_{A}}=\frac{V_{B}}{T_{B}}$，对外做功 $W=p(V_{B}-V_{A})$，根据热力学第一定律 $\Delta U=Q-W$，解得 $\Delta U=5.0\times10^{2}\UJ$。

\textbf{B．}（1）旋转偏振片 P，A 处得到的是始终强度相同的偏振光，偏振光再经过偏振片，在 B 处的光强随着 P 转动而变化，当 Q 的透振方向与经过 P 的偏振光的振动方向垂直时，B 处的光强为 0。
（2）插在 D 点的大头针必须挡住 C 及 A、B 的像，这样才保证沿 A、B 的光线经过 C、D；作出光路图如图所示，以入射点 O 为圆心作圆，交入射光线与折射光线于 E、F，从 E、F 作法线的垂线交法线于 G、H，用刻度尺量出 EG、FH 的长，由公式 $n=\frac{\sin i}{\sin r}=\frac{\frac{EG}{EO}}{\frac{FH}{FO}}=\frac{EG}{FH}$ 求出折射率。
（3）设 P 波的传播时间为 $t$，则 $x=v_{P}t$，$x=v_{S}(t+\Delta t)$，解得 $x=\frac{v_{P}v_{S}}{v_{P}-v_{S}}\Delta t$。代入数据得 $x=1260\Ukm$。由 $\lambda=v_{S}T$，解得 $\lambda=4.2\Ukm$。

\textbf{C．}（1）从能级图上可以看出，a 光子的能量最大，a 光的波长最短，b 光子的能量最小，b 光的波长最长，因此 C 选项正确。
（2）核反应过程遵循质量数、电荷数守恒，因此 \ce{^{1}_{0}n + ^{1}_{1}H -> ^{2}_{1}H}；比结合能即平均每个核子释放的能量，两个核子结合放出 $Q$ 的能量，比结合能为 $\frac{Q}{2}$。
（3）光子能量 $\varepsilon=h\nu$，动量 $p=\frac{h}{\lambda}$，且 $\nu=\frac{c}{\lambda}$，得 $p=\frac{\varepsilon}{c}$，则 $p_{A}:p_{B}=2:1$。A 照射时，光电子的最大初动能 $E_{A}=\varepsilon_{A}-W_{0}$。同理，$E_{B}=\varepsilon_{B}-W_{0}$。解得 $W_{0}=E_{A}-2E_{B}$。
}

\item
%% number: 13
%% paperName: 2012年江苏高考第 13 题 15 分
%% typeId: 6
%% type: 计算题
%% chapter: 电磁感应
%% point: 切割磁感线电动势
%% method: 无
%% score: 15
%% degree: 450
%% duplicateId: 0
%% body:
某兴趣小组设计了一种发电装置，如图所示。在磁极和圆柱状铁芯之间形成的两磁场区域的圆心角 $\alpha$ 均为 $\frac{4}{9}\pi$，磁场均沿半径方向．匝数为 $N$ 的矩形线圈 $abcd$ 的边长 ab$=$cd$=l$、bc$=$ad$=2l$。线圈以角速度 $\omega$ 绕中心轴匀速转动，bc和ad边同时进入磁场。在磁场中，两条边所经过处的磁感应强度大小均为 $B$、方向始终与两边的运动方向垂直。线圈的总电阻为 $r$，外接电阻为 $R$。求：
\begin{enumerate}
	\item
	线圈切割磁感线时，感应电动势的大小 $E_{m}$；
	\item
	线圈切割磁感线时，bc边所受安培力的大小 $F$；
	\item
	外接电阻上电流的有效值 $I$。
\end{enumerate}
\onepicture[w=6.0cm, align=right]{figs/13.png}
%% answer: 
\jdanswer{
\begin{enumerate}
	\item
	$E_{m}=2NBl^{2}\omega$
	\item
	$F=\frac{4N^{2}B^{2}l^{2}\omega}{r+R}$
	\item
	$I=\frac{4N^{2}B^{2}l^{2}\omega}{3(r+R)}$
\end{enumerate}
}
%% memo:
\memoanswer{
（1）bc、ad 边的运动速度 $v=\omega\frac{l}{2}$，感应电动势 $E_{m}=4NBlv$，解得 $E_{m}=2NBl^{2}\omega$。
（2）电流 $I_{m}=\frac{E_{m}}{r+R}$，安培力 $F=2NBI_{m}l$，解得 $F=\frac{4N^{2}B^{2}l^{2}\omega}{r+R}$。
（3）一个周期内，通电时间 $t=\frac{4}{9}T$，R 上消耗的电能 $W=I_{m}^{2}Rt$ 且 $W=I^{2}RT$，解得 $I=\frac{4N^{2}B^{2}l^{2}\omega}{3(r+R)}$。
}

\item
%% number: 14
%% paperName: 2012年江苏高考第 14 题 16 分
%% typeId: 6
%% type: 计算题
%% chapter: 功和能
%% point: 动能定理
%% method: 分情况讨论
%% score: 16
%% degree: 450
%% duplicateId: 0
%% body:
某缓冲装置的理想模型如图所示，劲度系数足够大的轻质弹簧与轻杆相连，轻杆可在固定的槽内移动，与槽间的滑动摩擦力恒为 $f$。轻杆向右移动不超过 $l$ 时，装置可安全工作。一质量为 $m$ 的小车若以速度 $v_{0}$ 撞击弹簧，将导致轻杆向右移动 $\frac{l}{4}$。轻杆与槽间的最大静摩擦力等于滑动摩擦力，且不计小车与地面的摩擦。
\begin{enumerate}
	\item
	若弹簧的劲度系数为 $k$，求轻杆开始移动时，弹簧的压缩量 $x$；
	\item
	求为使装置安全工作，允许该小车撞击的最大速度 $v_{m}$；
	\item
	讨论在装置安全工作时，该小车弹回速度 $v'$ 和撞击速度 $v$ 的关系。
\end{enumerate}
\onepicture[w=6.4cm, align=right]{figs/14.png}
%% answer: 
\jdanswer{
\begin{enumerate}
	\item
	$\frac{f}{k}$
	\item
	$\sqrt{v_{0}^{2}+\frac{3fl}{2m}}$
	\item
	当 $v<\sqrt{v_{0}^{2}-\frac{fl}{2m}}$ 时，$v'=v$；当 $\sqrt{v_{0}^{2}-\frac{fl}{2m}}<v<\sqrt{v_{0}^{2}+\frac{3fl}{2m}}$ 时，$v'=\sqrt{v_{0}^{2}-\frac{fl}{2m}}$。
\end{enumerate}
}
%% memo:
\memoanswer{
（1）轻杆开始移动时，弹簧的弹力 $F=kx$，且 $F=f$，解得 $x=\frac{f}{k}$。
（2）设轻杆移动前小车对弹簧所做的功为 $W$，则小车从撞击到停止的过程中，由动能定理得 $-f\cdot\frac{l}{4}-W=0-\frac{1}{2}mv_{0}^{2}$；同理，小车以 $v_{m}$ 撞击弹簧时，$-fl-W=0-\frac{1}{2}mv_{m}^{2}$，解得 $v_{m}=\sqrt{v_{0}^{2}+\frac{3fl}{2m}}$。
（3）设轻杆恰好移动时，小车撞击速度为 $v_{1}$，则 $\frac{1}{2}mv_{1}^{2}=W$，由上述两式解得 $v_{1}=\sqrt{v_{0}^{2}-\frac{fl}{2m}}$。当 $v<\sqrt{v_{0}^{2}-\frac{fl}{2m}}$ 时，$v'=v$；当 $\sqrt{v_{0}^{2}-\frac{fl}{2m}}<v<\sqrt{v_{0}^{2}+\frac{3fl}{2m}}$ 时，$v'=\sqrt{v_{0}^{2}-\frac{fl}{2m}}$。
}

\item
%% number: 15
%% paperName: 2012年江苏高考第 15 题 16 分
%% typeId: 6
%% type: 计算题
%% chapter: 磁场
%% point: 带电粒子在复合场中的运动
%% method: 无
%% score: 16
%% degree: 250
%% duplicateId: 0
%% body:
如图所示，待测区域中存在匀强电场和匀强磁场，根据带电粒子射入时的受力情况可推测其电场和磁场。图中装置由加速器和平移器组成，平移器由两对水平放置、相距为 $l$ 的相同平行金属板构成，极板长度为 $l$、间距为 $d$，两对极板间偏转电压大小相等、电场方向相反．质量为 $m$、电荷量为 $+q$ 的粒子经加速电压 $U_{m}$ 加速后，水平射入偏转电压为 $U_{1}$ 的平移器，最终从A点水平射入待测区域。不考虑粒子受到的重力。

\onepicture[w=9.0cm]{figs/15.png}

\begin{enumerate}
	\item
	求粒子射出平移器时的速度大小 $v_{1}$；
	\item
	当加速电压变为 4$U_{0}$ 时，欲使粒子仍从A点射入待测区域，求此时的偏转电压 $U$；
	\item
	已知粒子以不同速度水平向右射入待测区域，刚进入时的受力大小均为 $F$。现取水平向右为 $x$ 轴正方向，建立如图所示的直角坐标系 $Oxyz$。保持加速电压为 $U_{0}$ 不变，移动装置使粒子沿不同的坐标轴方向射入待测区域，粒子刚射入时的受力大小如下表所示。
	
	\begin{center}
	\begin{tabular}{|c|c|c|c|c|}
	\hline
	射入方向 & $y$ & $-y$ & $z$ & $-z$ \\
	\hline
	受力大小 & $\sqrt{5}F$ & $\sqrt{5}F$ & $\sqrt{7}F$ & $\sqrt{3}F$ \\
	\hline
	\end{tabular}
	\end{center}
	
	请推测该区域中电场强度和磁感应强度的大小及可能的方向。
\end{enumerate}
%% answer: 
\jdanswer{
\begin{enumerate}
	\item
	$\sqrt{\frac{2qU_{0}}{m}}$
	\item
	4$U_{1}$
	\item
	$E=\frac{F}{q}$，$B$ 平行于 $x$ 轴；$B=\frac{F}{q}\sqrt{\frac{2m}{qU_{0}}}$；$E$ 与 $Oxy$ 平面平行且与 $x$ 轴正方向的夹角为 30${^\circ}$ 或 150${^\circ}$（$B$ 沿 $x$ 轴方向），或 $-30{^\circ}$ 或 $-150{^\circ}$（$B$ 沿 $-x$ 轴方向）。
\end{enumerate}
}
%% memo:
\memoanswer{
（1）设粒子射出加速器的速度为 $v_{0}$，由动能定理 $qU_{0}=\frac{1}{2}mv_{0}^{2}$，由题意得 $v_{1}=v_{0}$，即 $v_{1}=\sqrt{\frac{2qU_{0}}{m}}$。
（2）在第一个偏转电场中，设粒子的运动时间为 $t$，加速度的大小 $a=\frac{qU_{1}}{md}$；在离开时，竖直分速度 $v_{y}=at$，竖直位移 $y_{1}=\frac{1}{2}at^{2}$，水平位移 $l=v_{1}t$。粒子在两偏转电场间做匀速直线运动，经历时间也为 $t$，竖直位移 $y_{2}=v_{y}t$。由题意知，粒子竖直总位移 $y=2y_{1}+y_{2}$，解得 $y=\frac{U_{1}l^{2}}{U_{0}d}$。则当加速电压为 4$U_{0}$ 时，$U=4U_{1}$。
（3）（a）由沿 $x$ 轴方向射入时的受力情况可知：$B$ 平行于 $x$ 轴，且 $E=\frac{F}{q}$。
（b）由沿 $\pm y$ 轴方向射入时的受力情况可知：$E$ 与 $Oxy$ 平面平行。由 $F^{2}+f^{2}=(\sqrt{5}F)^{2}$，则 $f=F$ 且 $f=qv_{1}B$，解得 $B=\frac{F}{q}\sqrt{\frac{2m}{qU_{0}}}$。
（c）设电场方向与 $x$ 轴方向夹角为 $\alpha$。若 $B$ 沿 $x$ 轴方向，由沿 $z$ 轴方向射入时的受力情况得 $(f+F\sin\alpha)^{2}+(F\cos\alpha)^{2}=(\sqrt{7}F)^{2}$，解得 $\alpha=30^{\circ}$ 或 $150^{\circ}$，即 $E$ 与 $Oxy$ 平面平行且与 $x$ 轴方向的夹角为 30${^\circ}$ 或 150${^\circ}$。同理，若 $B$ 沿 $-x$ 轴方向，$E$ 与 $Oxy$ 平面平行且与 $x$ 轴方向的夹角为 $-30{^\circ}$ 或 $-150{^\circ}$。
}

\end{enumerate}

\end{document}
